\begin{figure}[tbp]
\centering
\begin{tikzpicture}[>=Stealth,font=\small,x=0.75cm,y=1cm,
    every node/.style={inner sep=1.5pt},
    cell/.style={draw,minimum height=0.5cm,inner sep=0pt},
    lb/.style={font=\scriptsize,inner sep=1pt},
    side/.style={font=\footnotesize,anchor=west,align=left},
    lvl/.style={font=\footnotesize,anchor=east}]
  % level 0: w = 010#10 with its run
  \foreach \l/\f [count=\i from 0] in {0/white,1/black!12,0/black!12,\sep/white,1/black!12,0/black!12}
    \draw[fill=\f] (\i,-0.25) rectangle node{$\l$} (\i+1,0.25);
  \foreach \q [count=\i from 0] in {0,3,0,3,4,2,0}
    \node[lb,anchor=north] at (\i,-0.28) {$\q$};
  \draw[very thick,->] (1,-0.6) .. controls (1.6,-1.25) and (3.4,-1.25) .. (4,-0.6);
  \node[lvl] at (-0.3,0) {level $0$: $w$};
  \node[lvl,font=\scriptsize] at (-0.3,-0.42) {run};
  \node[side] at (6.5,-0.45) {(R): the anchored comparison $(1,4)$,\\
    state pair $(3,4)\in\Pi$, continuation letter $\sep$};
  % level 1: w_1 = abc, marker 0
  \draw (0,1.05) rectangle node{$a$} (1,1.55);
  \draw (1,1.05) rectangle node{$b=[1]$} (3,1.55);
  \draw (3,1.05) rectangle node{$c=[\sep1]$} (6,1.55);
  \node[lb,anchor=south] at (0.5,1.55) {$[\eps]$};
  \foreach \x in {0,1,3,6} \draw[densely dotted] (\x,0.25) -- (\x,1.05);
  \draw[very thick] (3,1.05) -- (3,1.55);
  \node[lvl] at (-0.3,1.3) {level $1$: $w_1$};
  \node[side] at (6.5,1.3) {(S$'$) at the cut $2$ of $w_1$, the terminal comparison\\
    $(2,3)$, continuation letter $c\in\Phi_1$};
  % level 2: w_2 = [ab], marker c
  \draw (0,2.1) rectangle node{$[ab]$} (6,2.6);
  \foreach \x in {0,6} \draw[densely dotted] (\x,1.55) -- (\x,2.1);
  \node[lvl] at (-0.3,2.35) {level $2$: $w_2$};
  \node[side] at (6.5,2.35) {(F): the positive terminal test at the\\
    position $\len{ab}$ of the block, letter $c\in\Phi_1$};
  % after deleting the column of b
  \begin{scope}[yshift=-2.35cm]
    \foreach \l/\f [count=\i from 0] in {0/white,\sep/white,1/white,0/white}
      \draw[fill=\f] (\i,-0.25) rectangle node{$\l$} (\i+1,0.25);
    \foreach \q [count=\i from 0] in {0,3,4,2,0}
      \node[lb,anchor=north] at (\i,-0.28) {$\q$};
    \draw[densely dashed,->] (0,0.3) .. controls (0.6,0.85) and (2.4,0.85) .. (3,0.3);
    \draw[densely dashed,->] (1,0.3) .. controls (1.6,0.85) and (3.4,0.85) .. (4,0.3);
    \node[lvl] at (-0.3,0) {without $b$};
    \node[side] at (6.5,0.05) {anchored comparisons $(0,3)$ and $(1,4)$,\\
      state pairs $(0,2)$ and $(3,0)$, both outside $\Pi$};
  \end{scope}
\end{tikzpicture}
\caption{A positive witness travels up the hierarchy
(\cref{ex:mo-loses-tower,ex:mo-loses}). Bottom up: the hole $w=010\sep10$ of
the reader of \cref{ex:mo-loses-tower} with its run, the word $w_1=abc$ of
blocks for the marker $0$, each letter drawn over its block and the marker that
closes it, and the one-letter word $w_2=[ab]$ for the marker $c$. The anchored
witness $(1,4)$ in $\Pi$, whose matched factor $10$ is shaded, is clause (R) of
the base channel at level $0$, clause (S$'$) of the derived channel at level
$1$, and clause (F) at level $2$, where it is a test of the final label
(\cref{def:mo-channel,lem:mo-desub}). The compression for plain entailment sees
no difference between the prefixes $a$ and $ab$ of the block and deletes the
column of $b$, which carries the matched factor $10$; the result $0\sep10$
(last row) is still a hole, but no longer a $\Pi$-hole. Recording the channel in
the summaries keeps $b$.}
\Description{Three stacked rows of cells. The bottom row is the word 0 1 0
separator 1 0, with the states 0 3 0 3 4 2 0 of its run under the cuts and a
thick arrow below from the cut 1 to the cut 4; the cells 1 0 before the
separator and the last cells 1 0 are shaded. The middle row has the letters a
over the first cell, b equal to the block 1 over the next two cells, and c
equal to the block separator 1 over the last three cells; the boundary between
b and c, the cut 2, is drawn thick. The top row is one cell, the block a b,
over the whole word. Notes to the right name the clause of the witness at each
level: R at level 0, S prime at level 1, F at level 2. A separate row at the
bottom shows the word 0 separator 1 0 with the run 0 3 4 2 0 and two dashed
arrows, from the cut 0 to the cut 3 and from the cut 1 to the cut 4, whose
state pairs are outside Pi.}
\label{fig:mo-tower}
\end{figure}
